\scanpage{61}
\mspage{67}

Donc la flèche envisagée ($**$) s'interprète comme

\[ H_\bullet^{\unsure{\mathbb D|\underline D}}({}_{y\backslash}X,M_X)\longrightarrow M \]

et c'est l'hom d'adjonction $p_!(p^*M))\unsure{\relbar\joinrel\relbar|\longrightarrow}M$. Donc c'est un iso pour tt $M$ ssi $ {}_{y\backslash}X$ est \deleted{\unsure{$H_{\mathbb D}$}} $\mathbb D$-contractile.

On va \deleted{\unsure{ré}} énoncer ce qu'on a prouvé

\uline{Théorème 1} Soit

\diagram{\begin{tikzcd}[ampersand replacement=\&,row sep=tiny,column sep=small]
f:X \arrow[rr,"f"] \arrow[dr,dash,"p"'] \& {} \& Y \arrow[dl,"q"] \\
{} \& e \& {}
\end{tikzcd}}

dans $\mathrm{Diag}$. On suppose que $\mathbb D$ est un \deleted{\unsure{dér}} prédérivateur satisfaisant Der $(2)_{\unsure{x|X|*}}$ \deleted{et Der \unsure{(4)(d)} pour $f$} que \deleted{\unsure{$f_*$}} $p_*,q_*$ et $f_!$ existent, que $f_!$ satisfait à Der 4 (g) pour le calcul des fibres. Conditions équivalentes

a) $f$ est une $\mathbb D$-équiv. costricte, i.e. (prop. 2) $\forall y\in\mathrm{Ob}\,Y$, $ {}_{y\backslash}X$ est $\mathbb D$-contractile

b) $\forall F\in\mathrm{Ob}\,\mathbb D(Y)$, l'hom suivant est iso

\[ H^\bullet_{\mathbb D}(Y,F)\mathrel{\overset{\sim}{\unsure{\relbar\joinrel\relbar|\longrightarrow}}}H^\bullet_{\mathbb D}(X,f^*F) \]

c) $\forall M\in\mathrm{Ob}(A=\mathbb D(e))$, l'hom d'adjonction est iso.

\[ \alpha_!\underbrace{\alpha^*(M_Y)}_{M_X}\xrightarrow{\sim}M_Y. \]

Je vais renuméroter les \unsure{«|4} axiomes sur les prédérivateurs,

Der 1, Dér 2 comme avant (p. 45)

% slanted
\margin{Axiome des \deleted{diag} cat. d'indices asphériques}

\uline{Dér 3} : Soit $X$ dans $\mathrm{Diag}$, $p:X\mathrel{\relbar\joinrel\relbar}e$. \inserted{Supposons $X$ \uline{contractile}} Alors

\[ p^*:M\longmapsto M_X:\underset{\substack{\Vert\\\mathbb D(e)}}{A}\longrightarrow\mathbb D(X) \]

est pleinement fidèle.

\margin{\deleted{\unsure{En}} \unsure{présence|préciser} \unsure{de|des} Dér (3,4) NB Il suffit de supposer que $X$ a un objet \uline{final}. Mais ce n'est \deleted{\illegible{3}}}

% slanted
\margin{pas un axiome \unsure{autodual|auto-dual} !}

\margin{$\searrow$ cf \uline{variantes} Der a) b) c) page 94, et \unsure{suivantes|surtout} p$^{\mathrm s}$ \deleted{\unsure{96}} \unsure{111--112|111,112}}
% archive sheet number: 60
